\chapter{Conclusions}

\section{Summary of the Study}

This thesis investigated whether a small, open-source large language model could be adapted for the task of first-innings score prediction in T20 cricket through parameter-efficient fine-tuning on domain-specific data. The central argument was that the structured, context-rich nature of in-game cricket match states could be encoded as natural-language prompts, enabling a fine-tuned LLM to learn the mapping from mid-innings conditions to final scores without the need for an explicit probabilistic model of ball outcomes.

To test this argument, a complete pipeline was constructed: raw ball-by-ball match data from the Cricsheet archive was preprocessed and enriched with engineered features, including a venue par score derived from historical ground averages, a Role-Based Scoring System (RBSS) that quantifies remaining batting depth as a weighted sum of player role contributions, and player classifications for active batters and bowlers. These features were serialised into a structured natural-language prompt schema, producing one training example per over boundary per innings. After temporal filtering, 6,095 training examples covering approximately 2,895 international men's T20 matches were used to fine-tune \textit{Llama 3.2-3B} using QLoRA (4-bit NF4 quantisation with rank-8 LoRA adapters) on a Kaggle T4 GPU. The fine-tuned model was then evaluated on a held-out temporal test set alongside two zero-shot baselines: the unmodified \textit{Llama 3.2-3B} base model and the frontier model \textit{GPT-4.1-nano}.

\section{Principal Findings}

The experiments produced three principal findings.

\textbf{First}, domain-specific fine-tuning transforms an LLM with no useful prior knowledge for this task into a competitive score predictor. The unmodified \textit{Llama 3.2-3B} model achieves an MAE of 71.31 runs and an $R^2$ of $-$317.4\% in the zero-shot setting, confirming that general pre-training provides essentially no usable knowledge for structured numerical prediction from mid-innings cricket states. After fine-tuning on 6,095 domain-specific examples, the same model achieves an MAE of 16.51 runs and an $R^2$ of 68.3\% --- a 76.8\% reduction in MAE and a 385.7 percentage-point improvement in $R^2$.

\textbf{Second}, a small fine-tuned model substantially outperforms a much larger frontier model in the zero-shot setting. \textit{GPT-4.1-nano}, with approximately 100 times the parameter count of \textit{Llama 3.2-3B}, achieves an MAE of 30.02 runs and an $R^2$ of 14.0\% without fine-tuning --- considerably worse than the fine-tuned model on both metrics. This result demonstrates that scale alone does not substitute for domain-specific supervised training on structured quantitative data, and that parameter-efficient fine-tuning of a small open-source model is a viable and cost-effective alternative to prompt-engineering a large proprietary frontier model for specialised prediction tasks.

\textbf{Third}, checkpoint selection based on minimum validation loss provides a marginal but consistent generalisation benefit. The step-800 checkpoint, which corresponds to the minimum evaluation loss during training, outperforms the final epoch checkpoint on all three metrics when evaluated on 500 test datapoints (MAE: 16.51 vs.\ 16.54 runs; $R^2$: 68.3\% vs.\ 68.1\%), confirming the standard recommendation to prefer minimum-validation-loss checkpoints over final-epoch checkpoints in fine-tuning workflows.

\section{Contributions}

This study makes three primary contributions to the fields of sports analytics and applied natural language processing.

\textbf{Contribution 1 --- A complete, reproducible T20 score prediction pipeline.} The end-to-end pipeline from raw Cricsheet JSON data through feature engineering, prompt generation, QLoRA fine-tuning, and evaluation is fully documented and versioned, providing a reproducible baseline for future work in LLM-based cricket analytics.

\textbf{Contribution 2 --- The Role-Based Scoring System (RBSS) feature.} The RBSS is a novel engineered feature that quantifies the expected run-scoring potential of the batting depth remaining in an innings, weighted by the assigned role of each player. This feature provides the model with information about the batting lineup that is not recoverable from the current score and wickets-in-hand alone, and its inclusion in the prompt schema is a key factor in enabling the model to learn over-boundary score dynamics.

\textbf{Contribution 3 --- Empirical evidence for fine-tuned LLM superiority over frontier zero-shot models on domain-specific numerical prediction.} The comparison between fine-tuned \textit{Llama 3.2-3B} and zero-shot \textit{GPT-4.1-nano} provides a concrete, quantitative illustration of a principle that has been observed in the broader LLM literature but has not previously been demonstrated in a cricket analytics context: that domain-specific fine-tuning of a small open-source model yields substantially better results than zero-shot prompting of a large frontier model for structured numerical prediction tasks.

\section{Recommendations for Practitioners}

Based on the findings of this study, the following practical recommendations are offered to data scientists and machine learning engineers working on sports prediction problems.

When the prediction target is a structured numerical outcome derived from a well-defined in-game state (such as a cricket final total, a basketball final score, or a football goal tally), feature engineering to encode domain knowledge in the prompt is at least as important as model selection. A well-engineered prompt schema that includes venue context, remaining batting resources, and opponent quality will enable a small fine-tuned model to outperform a frontier model receiving only raw in-game statistics.

For prediction tasks involving high stochasticity (e.g., early-innings score prediction where fewer than 5 overs have been bowled), point prediction from a single model completion is likely to produce poorly calibrated estimates. Practitioners are encouraged to report prediction intervals or to evaluate multiple sampled completions in order to quantify uncertainty.

For fine-tuning workflows with limited GPU memory, QLoRA with rank-8 adapters and 4-bit NF4 quantisation provides an effective trade-off between model quality and memory footprint, and the minimum-validation-loss checkpoint should be adopted in preference to the final epoch checkpoint.

\section{Concluding Remarks}

This thesis demonstrates that context-aware fine-tuning of large language models on domain-specific, structured natural-language prompts is a viable and effective approach for first-innings score prediction in T20 cricket. The work makes a case for the broader applicability of parameter-efficient fine-tuning as a practical methodology for sports analytics: one that leverages the representational power of pre-trained language models while remaining tractable on consumer-grade hardware. The pipeline, features, and experimental framework developed here provide a foundation on which future work can build --- extending to domestic leagues, larger models, second-innings prediction, and live real-time deployment.
